\section{Armesa II}


\subsection{Presentación del caso}

La Dirección Comercial de ARMESA ha conseguido nuevos contratos con clientes mayoristas,
eliminando las ventas directas a minoristas. Las entregas mensuales pasan a ser
\emph{exactas} y se recogen en la Tabla~\ref{tab:armesaII-entregas}.

\begin{table}[!h]
\centering
\caption{entregas mensuales exigidas por producto.}
\label{tab:armesaII-entregas}
\begin{tabular}{lcccc}
\hline
Producto & 1 & 2 & 3 & 4 \\ \hline
Unidades & 3000 & 500 & 1000 & 2000 \\
\hline
\end{tabular}
\end{table}

Se han encontrado talleres externos a los que puede subcontratarse la \textbf{embutición} y el
\textbf{mecanizado} de cualquiera de los productos, lo que incrementa en un 20\% los costes
variables unitarios de las unidades subcontratadas. El mismo proveedor realiza ambas
operaciones y adquiere la chapa; por ello, la limitación de 2\,000 pies cuadrados mensuales
de chapa especial para los productos 2 y 4 sólo afecta a la producción \emph{interna}
(embutida y mecanizada en ARMESA):contentReference[oaicite:1]{index=1}.

La empresa puede disponer de hasta 100 horas extraordinarias mensuales en el taller de
\textbf{acabado}. Cuando se utilizan, encarecen los costes variables unitarios según
la Tabla~\ref{tab:armesaII-recargos}.

\begin{table}[!h]
\centering
\caption{recargo unitario en costes por uso de horas extra de acabado (€/unidad).}
\label{tab:armesaII-recargos}
\begin{tabular}{lcccc}
\hline
Producto & 1 & 2 & 3 & 4 \\ \hline
Recargo  & 0{,}20 & 0{,}40 & 0{,}30 & 0{,}30 \\
\hline
\end{tabular}
\end{table}

El objetivo sigue siendo formular un modelo de programación lineal que determine el
programa mensual de producción óptimo bajo estas nuevas condiciones.


\subsection{Formulación explícita}

\paragraph{Variables}\mbox{}\\

Sean $x_{ijk}$ las unidades producidas, donde:
\begin{itemize}
    \item $i$: tipo de producto ($i=1,2,3,4$),
    \item $j$: lugar de embutición + mecanizado ($j=0$: interno, $j=1$: subcontratado),
    \item $k$: tipo de acabado ($k=0$: normal, $k=1$: horas extra).
\end{itemize}

\paragraph{Restricciones}\mbox{}\\
\\
Disponibilidad de embutición:
\begin{align}
0.03(x_{100}+x_{101}) + 0.15(x_{200}+x_{201}) & \notag \\
+ 0.05(x_{300}+x_{301}) + 0.10(x_{400}+x_{401}) & \leq 400
\end{align}

Disponibilidad de mecanizado:
\begin{align}
0.06(x_{100}+x_{101}) & \notag \\
+ 0.12(x_{200}+x_{201}) & \notag \\
+ 0.10(x_{400}+x_{401}) & \leq 400
\end{align}


Disponibilidad de montaje:
\begin{align}
0.02(x_{100}+x_{101}+x_{110}+x_{111}) & \notag \\
+0.06(x_{200}+x_{201}+x_{210}+x_{211}) & \notag \\
+0.02(x_{300}+x_{301}+x_{310}+x_{311}) & \notag \\
+0.05(x_{400}+x_{401}+x_{410}+x_{411}) & \leq 500 \label{eq:armesa-montaje}
\end{align}


Disponibilidad de acabado normal:
\begin{align}
0.04x_{100}+0.20x_{200}+0.03x_{300}+0.12x_{400} & \notag \\
+0.04x_{110}+0.20x_{210}+0.03x_{310}+0.12x_{410} & \leq 450
\end{align}


Disponibilidad de horas extra de acabado:
\begin{align}
0.04x_{101}+0.20x_{201}+0.03x_{301}+0.12x_{401} & \notag \\
+0.04x_{111}+0.20x_{211}+0.03x_{311}+0.12x_{411} & \leq 100
\end{align}


Disponibilidad de embalaje:
\begin{align}
0.02(x_{100}+x_{101}+x_{110}+x_{111}) & \notag \\
+0.06(x_{200}+x_{201}+x_{210}+x_{211}) & \notag \\
+0.02(x_{300}+x_{301}+x_{310}+x_{311}) & \notag \\
+0.05(x_{400}+x_{401}+x_{410}+x_{411}) & \leq 400
\label{eq:armesa-embalaje}
\end{align}

Disponibilidad de chapa:
\begin{equation}
2(x_{200}+x_{201}) + 1.2(x_{400}+x_{401}) \leq 2000
\end{equation}

Restricciones de mercado (entregas exactas):
\begin{align}
x_{100}+x_{101}+x_{110}+x_{111} &= 3000 \notag\\
x_{200}+x_{201}+x_{210}+x_{211} &= 500 \notag \\
x_{300}+x_{301}+x_{310}+x_{311} &= 1000 \notag\\
x_{400}+x_{401}+x_{410}+x_{411} &= 2000
\end{align}


Conviene notar que, dado que las cantidades que hay que producir de cada producto son fijas y que el montaje y el embalaje se hace con las 500 y 400 horas disponbiles, respectivamente, las restricciones \ref{eq:armesa-montaje} y \ref{eq:armesa-embalaje} no dependen de los valores de las variables porque tanto el primer miembro de la desigualdad como el segundo son constantes. Dicho de otro modo, dado que el montaje y el emabalaje se hace con una horas determinadas y las cantidades están fijadas, o hay posibilidad de embalar y montar todas las unidades o no. Si no hay suficientes horas de uno o de otro proceso, no habría forma de obtener un plan de producción (el problema sería no factible).



\paragraph{Función objetivo}\mbox{}\\
\\
La función objetivo incorpora las contribuciones unitarias (precio de venta menos coste variable), 
penalizando subcontratación (+20\% en costes) y recargos por horas extra:

\begin{align}
\max z = & \; 4.0x_{100} + 3.8x_{101} + 2.8x_{110} + 2.6x_{111} \notag \\
&+ 10.0x_{200} + 9.6x_{201} + 7.0x_{210} + 6.6x_{211} \notag \\
&+ 5.0x_{300} + 4.8x_{301} + 2.8x_{310} + 2.6x_{311} \notag \\
&+ 6.0x_{400} + 5.7x_{401} + 3.2x_{410} + 2.9x_{411}
\end{align}


% Formulación explícita Armesa II (alineada por columnas)

\paragraph{Modelo completo}\mbox{}\\

{\small
\begin{align}
\max\; z \;=\;
&\; 4.0x_{100} + 3.8x_{101} + 2.8x_{110} + 2.6x_{111} \notag\\
&+ 10.0x_{200} + 9.6x_{201} + 7.0x_{210} + 6.6x_{211} \notag\\
&+ 5.0x_{300} + 4.8x_{301} + 6.0x_{400} + 5.7x_{401} \notag\\
&+ 3.2x_{410} + 2.9x_{411}
\\[4pt]
\text{s.a.}\quad
&0.03x_{100}+0.03x_{101}+0.15x_{200}+0.15x_{201}+0.05x_{300}+0.05x_{301} \notag\\
&+0.10x_{400}+0.10x_{401}\le400 \\[2pt]
&0.06x_{100}+0.06x_{101}+0.12x_{200}+0.12x_{201}+0.10x_{400}+0.10x_{401}\le400 \\[2pt]
&0.02(x_{100}+x_{101}+x_{110}+x_{111})
+0.06(x_{200}+x_{201}+x_{210}+x_{211}) \notag\\
&+0.02(x_{300}+x_{301}+x_{310}+x_{311})
+0.05(x_{400}+x_{401}+x_{410}+x_{411})\le500 \\[2pt]
&0.04x_{100}+0.04x_{110}+0.20x_{200}+0.20x_{210}
+0.03x_{300}+0.03x_{310}+0.12x_{400}+0.12x_{410}\leq \notag\ \\ & 450  \\[2pt]
&0.04x_{101}+0.04x_{111}+0.20x_{201}+0.20x_{211}
+0.03x_{301}+0.03x_{311}+0.12x_{401}+0.12x_{411}\leq \notag\ \\ & 100  \\[2pt]
&0.02(x_{100}+x_{101}+x_{110}+x_{111})
+0.06(x_{200}+x_{201}+x_{210}+x_{211}) \notag\\
&+0.02(x_{300}+x_{301}+x_{310}+x_{311})
+0.05(x_{400}+x_{401}+x_{410}+x_{411})\le400 \\[2pt]
&2.0(x_{200}+x_{201})+1.2(x_{400}+x_{401})\le2000 \\[2pt]
&x_{100}+x_{101}+x_{110}+x_{111}=3000 \\
&x_{200}+x_{201}+x_{210}+x_{211}=500 \\
&x_{300}+x_{301}+x_{310}+x_{311}=1000 \\
&x_{400}+x_{401}+x_{410}+x_{411}=2000 \\[2pt]
&x_{ijk}\ge0 \qquad i=1,\dots,4;\; j,k\in\{0,1\}
\end{align}
}




\subsection{Formulación generalizada}

\textbf{Conjuntos}\\[0.2cm]

\begin{tabularx}{\textwidth}{l c X}
    $\mathcal{P}$ & : & conjunto de productos \\
    $\mathcal{J}$ & : & conjunto de opciones de embutición y mecanizado ($0$: interno, $1$: subcontratado) \\
    $\mathcal{K}$ & : & conjunto de opciones de acabado ($0$: normal, $1$: horas extra) \\
    $\mathcal{T}$ & : & conjunto de talleres ($t=$ embutición, mecanizado, montaje, acabado, embalaje) \\
\end{tabularx}
\\[0.4cm]

\textbf{Parámetros}\\[0.2cm]

\begin{tabularx}{\textwidth}{l c X}
    $A_{tp}$   & : & horas requeridas en el taller $t \in \mathcal{T}$ para producir una unidad del producto $p \in \mathcal{P}$ (internamente) \\
    $H_t$      & : & horas disponibles en el taller $t \in \mathcal{T}$ \\
    $CH_p$     & : & cantidad de chapa necesaria por unidad de producto $p \in \mathcal{P}$ (sólo si se fabrica internamente) \\
    $CHAP$     & : & cantidad máxima disponible de chapa especial \\
    $P_p$      & : & precio de venta unitario del producto $p \in \mathcal{P}$ \\
    $C_p$      & : & coste variable unitario del producto $p \in \mathcal{P}$ \\
    $\alpha$   & : & factor de incremento del coste en caso de subcontratación ($\alpha = 1.2$) \\
    $\beta_p$  & : & recargo en el coste variable por unidad de producto $p$ en caso de acabado en horas extra \\
    $D_p$      & : & demanda mensual exacta del producto $p \in \mathcal{P}$ \\
    $H^{\text{extra}}$ & : & número máximo de horas de acabado extraordinarias disponibles \\
\end{tabularx}
\\[0.4cm]

\textbf{Variables}\\[0.2cm]

\begin{tabularx}{\textwidth}{l c X}
    $x_{p j k}$ & : & número de unidades producidas del producto $p \in \mathcal{P}$, con $j \in \mathcal{J}$ (interno/subcontratado) y $k \in \mathcal{K}$ (acabado normal/horas extra). \\
\end{tabularx}

\paragraph{Restricciones}\mbox{}\\

\noindent
Disponibilidad de talleres internos (aquellos talleres cuyos procesos no se subcontratan y donde no se usan horas extra: montaje y embalaje):
\begin{equation}
\sum_{p \in \mathcal{P}} \sum_{j \in \mathcal{J}} \sum_{k \in \mathcal{K}} A_{tp}\, x_{p j k} \;\; \leq \;\; H_t \qquad \forall t \in \mathcal{T}
\end{equation}

Disponibilidad de talleres cuyos procesos se subcontratan (embutición y mecanizado):
\begin{equation}
\sum_{p \in \mathcal{P}} \sum_{k \in \mathcal{K}} A_{tp}\, x_{p 0 k} \;\; \leq \;\; H_t \qquad \forall t \in \mathcal{T}
\end{equation}

\noindent
Disponibilidad de horas normales en acabado:
\begin{equation}
\sum_{p \in \mathcal{P}} \sum_{j \in \mathcal{J}} x_{p j 0  }\, A_{\text{acab},p} \;\; \leq \;\; H_t
\end{equation}

\noindent
Disponibilidad de horas extra en acabado:
\begin{equation}
\sum_{p \in \mathcal{P}} \sum_{j \in \mathcal{J}} x_{p j 1}\, A_{\text{acab},p} \;\; \leq \;\; H^{\text{extra}}
\end{equation}

\noindent
Disponibilidad de chapa (solo fabricación interna):
\begin{equation}
\sum_{p \in \mathcal{P}} \sum_{k \in \mathcal{K}} CH_p \, x_{p0k} \;\; \leq \;\; CHAP
\end{equation}

\noindent
Restricciones de mercado (entregas exactas):
\begin{equation}
\sum_{j \in \mathcal{J}} \sum_{k \in \mathcal{K}} x_{p j k} = D_p \qquad \forall p \in \mathcal{P}
\end{equation}
\\



\paragraph{Función objetivo}\mbox{}\\

\begin{align}
\max z = 
& \sum_{p \in \mathcal{P}} \Bigg[ 
\sum_{k \in \mathcal{K}} (P_p - C_p)\, x_{p0k} 
+ \sum_{k \in \mathcal{K}} (P_p - \alpha C_p)\, x_{p1k} \notag \\
& \qquad - \sum_{j \in \mathcal{J}} \beta_p \, x_{p j 1} 
\Bigg]
\end{align}
\\

\paragraph{Modelo completo}\mbox{}\\
\begin{align}
\mbox{max.} \quad & \sum_{p \in \mathcal{P}} \Bigg[ 
\sum_{k \in \mathcal{K}} (P_p - C_p)\, x_{p0k} 
+ \sum_{k \in \mathcal{K}} (P_p - \alpha C_p)\, x_{p1k} 
- \sum_{j \in \mathcal{J}} \beta_p \, x_{p j 1} \Bigg] \notag \\
\mbox{s. a.:} \quad 
& \sum_{p \in \mathcal{P}} \sum_{j \in \mathcal{J}} \sum_{k \in \mathcal{K}} A_{tp}\, x_{p j k} \leq H_t \qquad \forall t \in \mathcal{T} \notag \\
& \sum_{p \in \mathcal{P}} \sum_{j \in \mathcal{J}} A_{\text{acab},p}\, x_{p j 1} \leq H^{\text{extra}} \notag \\
& \sum_{p \in \mathcal{P}} \sum_{k \in \mathcal{K}} CH_p \, x_{p0k} \leq CHAP \notag \\
& \sum_{j \in \mathcal{J}} \sum_{k \in \mathcal{K}} x_{p j k} = D_p \qquad \forall p \in \mathcal{P} \notag \\
& x_{p j k} \geq 0 \qquad \forall p \in \mathcal{P}, \; j \in \mathcal{J}, \; k \in \mathcal{K} \notag
\end{align}








\subsection{Formulación explícita ($x,y,z$)}

% ===================== VARIABLES =====================
\textbf{Variables}\\[0.2cm]
\begin{tabular}{l c l}
    $x_i$           & : & unidades totales del producto $i$, \; $i=1,2,3,4$ \\
    $y_i^{\mathrm{I}}$ & : & unidades del producto $i$ con embutición y mecanizado internos \\
    $y_i^{\mathrm{S}}$ & : & unidades del producto $i$ con embutición y mecanizado subcontratados \\
    $z_i^{\mathrm{N}}$ & : & unidades del producto $i$ acabadas en turno normal \\
    $z_i^{\mathrm{E}}$ & : & unidades del producto $i$ acabadas en horas extra \\
\end{tabular}

% ===================== ENTREGAS EXACTAS Y FLUJOS =====================
\textbf{Entregas exactas y vinculación de flujos}\\[0.2cm]
\begin{align}
& x_1 = 3000, \quad x_2 = 500, \quad x_3 = 1000, \quad x_4 = 2000 \\
& x_i = y_i^{\mathrm{I}} + y_i^{\mathrm{S}} && i=1,2,3,4 \\
& x_i = z_i^{\mathrm{N}} + z_i^{\mathrm{E}} && i=1,2,3,4
\end{align}

% ===================== CAPACIDADES INTERNAS (EMB Y MEC) =====================
\textbf{Capacidades internas de embutición y mecanizado}\\[0.2cm]
\begin{align}
& 0.03\,y_1^{\mathrm{I}} + 0.15\,y_2^{\mathrm{I}} + 0.05\,y_3^{\mathrm{I}} + 0.10\,y_4^{\mathrm{I}} \le 400
&& \text{(embutición)} \\
& 0.06\,y_1^{\mathrm{I}} + 0.12\,y_2^{\mathrm{I}} \;+\; 0.10\,y_4^{\mathrm{I}} \le 400
&& \text{(mecanizado)}
\end{align}

% ===================== CAPACIDADES QUE DEPENDEN DEL TOTAL =====================
\textbf{Montaje y embalaje}\\[0.2cm]
\begin{align}
& 0.02\,x_1 + 0.06\,x_2 + 0.02\,x_3 + 0.05\,x_4 \le 500
&& \text{(montaje)} \\
& 0.02\,x_1 + 0.06\,x_2 + 0.02\,x_3 + 0.05\,x_4 \le 400
&& \text{(embalaje)}
\end{align}

% ===================== ACABADO NORMAL Y HORAS EXTRA =====================
\textbf{Acabado (normal y horas extra)}\\[0.2cm]
\begin{align}
& 0.04\,z_1^{\mathrm{N}} + 0.20\,z_2^{\mathrm{N}} + 0.03\,z_3^{\mathrm{N}} + 0.12\,z_4^{\mathrm{N}} \le 450
&& \text{(acabado normal)} \\
& 0.04\,z_1^{\mathrm{E}} + 0.20\,z_2^{\mathrm{E}} + 0.03\,z_3^{\mathrm{E}} + 0.12\,z_4^{\mathrm{E}} \le 100
&& \text{(acabado horas extra)}
\end{align}

% ===================== CHAPA ESPECIAL =====================
\textbf{Disponibilidad de chapa especial (solo producción interna)}\\[0.2cm]
\begin{align}
& 2.0\,y_2^{\mathrm{I}} + 1.2\,y_4^{\mathrm{I}} \le 2000
\end{align}

% ===================== FUNCIÓN OBJETIVO =====================


% ===================== NO NEGATIVIDAD =====================
\textbf{No negatividad}\\[0.2cm]
\begin{align}
& x_i \ge 0,\quad y_i^{\mathrm{I}} \ge 0,\quad y_i^{\mathrm{S}} \ge 0,\quad
  z_i^{\mathrm{N}} \ge 0,\quad z_i^{\mathrm{E}} \ge 0
\qquad i=1,2,3,4
\end{align}


Dado que los ingresos y los costes diferentes de las horas extra y la subcontratación son fijos (las cantidades producidas son fijas) el problema consiste en minimizar los costes por subcontratación más los costes por el huso de horas extra:

\begin{align}
\mbox{min.} \; z \;=\;
&\; (0.2\cdot 6)\,y_1^{\mathrm{S}} + (0.2\cdot 15)\,y_2^{\mathrm{S}} + (0.2\cdot 11)\,y_3^{\mathrm{S}} + (0.2\cdot 14)\,y_4^{\mathrm{S}} \notag \\
&\;+\; 0.20\,z_1^{\mathrm{E}} + 0.40\,z_2^{\mathrm{E}} + 0.20\,z_3^{\mathrm{E}} + 0.30\,z_4^{\mathrm{E}}
\end{align}


\textbf{Modelo completo}
\begin{align}
\min\; z \;=\;
&\; (0.2\cdot 6)\,y_1^{\mathrm{S}} + (0.2\cdot 15)\,y_2^{\mathrm{S}}
   + (0.2\cdot 11)\,y_3^{\mathrm{S}} + (0.2\cdot 14)\,y_4^{\mathrm{S}} \notag \\
&\;+\; 0.20\,z_1^{\mathrm{E}} + 0.40\,z_2^{\mathrm{E}}
       + 0.20\,z_3^{\mathrm{E}} + 0.30\,z_4^{\mathrm{E}}
\\[4pt]
\text{s.a.}\quad
& x_1 = 3000 \notag \\ 
& x_2 = 500\notag \\
& x_3 = 1000\notag \\
& x_4 = 2000 \notag \\
& x_i = y_i^{\mathrm{I}} + y_i^{\mathrm{S}} \qquad i=1,2,3,4 \notag \\
& x_i = z_i^{\mathrm{N}} + z_i^{\mathrm{E}} \qquad i=1,2,3,4 \notag \\
& 0.03\,y_1^{\mathrm{I}} + 0.15\,y_2^{\mathrm{I}} + 0.05\,y_3^{\mathrm{I}} + 0.10\,y_4^{\mathrm{I}} \le 400 \notag \\
& 0.06\,y_1^{\mathrm{I}} + 0.12\,y_2^{\mathrm{I}} + 0.10\,y_4^{\mathrm{I}} \le 400 \notag \\
& 0.02\,x_1 + 0.06\,x_2 + 0.02\,x_3 + 0.05\,x_4 \le 500 \notag \\
& 0.02\,x_1 + 0.06\,x_2 + 0.02\,x_3 + 0.05\,x_4 \le 400 \notag \\
& 0.04\,z_1^{\mathrm{N}} + 0.20\,z_2^{\mathrm{N}} + 0.03\,z_3^{\mathrm{N}} + 0.12\,z_4^{\mathrm{N}} \le 450 \notag \\
& 0.04\,z_1^{\mathrm{E}} + 0.20\,z_2^{\mathrm{E}} + 0.03\,z_3^{\mathrm{E}} + 0.12\,z_4^{\mathrm{E}} \le 100 \notag \\
& 2.0\,y_2^{\mathrm{I}} + 1.2\,y_4^{\mathrm{I}} \le 2000 \notag \\
& x_i,\,y_i^{\mathrm{I}},\,y_i^{\mathrm{S}},\,z_i^{\mathrm{N}},\,z_i^{\mathrm{E}} \ge 0 \qquad i=1,2,3,4 \notag
\end{align}


